\subsection{A preference round}\label{sec:pref-round}

Whereas a warm-up round (Section~\ref{sec:warmup}) supervises the model
with the gold SQL query alone, a preference round trains it to prefer,
among its own samples, those that reproduce the execution result of the
gold SQL query on the client's database. In round $t$, with
$T_w < t \le T$, the server sends $\theta^{t-1}$ to all $K$ clients
(Algorithm~\ref{alg:method}).

For each $(q, d, y^\star) \in \mathcal{D}_k$, client $k$ samples $N$
queries from $\pi_{\theta^{t-1}}(\cdot \mid x)$ at temperature
\setting{pref.t}, without top-$p$ or top-$k$ truncation and with at most
\setting{pref.tokens} new tokens; we set $N$ to \setting{pref.samples}.
It then executes $y^\star$ once and every sample on $d$. A sample $y$ is
labeled correct if $\mathrm{exec}(y, d) = \mathrm{exec}(y^\star, d)$,
with execution results compared as in the execution accuracy of the
benchmark, and incorrect otherwise. An incorrect sample is either an
execution error, which fails to execute or exceeds the time limit, or an
incorrect result, which executes and returns a different execution
result. A sample that contains no SQL query is labeled incorrect; a
sample truncated at the token limit is not labeled and enters no
preference pair. The execution label thus requires the gold SQL query as
well as the database, and only the client, which stores $d$, can compute
it.

For each question with at least one correct and one incorrect sample,
the client constructs $M$ preference pairs $(y^{+}, y^{-})$, where
$y^{+}$ is drawn uniformly from the correct samples of the question and
$y^{-}$ uniformly from its incorrect samples; we set $M$ to
\setting{pref.pairs}. The draw is over the samples rather than over the
distinct queries, and it is without replacement when the question has at
least $M$ samples of the label and with replacement otherwise. Since
both elements of a pair are samples of $\pi_{\theta^{t-1}}$, the
training signal targets errors that the current global model makes.

A question yields no pair when all its samples carry the same label,
when the tables of $d$ contain no rows, in which case the comparison of
execution results is uninformative, or when $y^\star$ fails to execute,
which precludes labeling. In the first case, the question is sampled
again in the next preference round.

Local training is initialized at $\theta^{t-1}$ and minimizes a loss
with two terms: a preference term, which increases the log-likelihood
ratio of the correct over the incorrect sample relative to a reference
model $\pi_{\mathrm{ref}}$, and a negative log-likelihood (NLL) term on
the correct sample. For a pair $(x, y^{+}, y^{-})$, we define
\begin{equation}
  m_\theta = \beta \left[
    \log \frac{\pi_\theta(y^{+} \mid x)}{\pi_{\mathrm{ref}}(y^{+} \mid x)}
    - \log \frac{\pi_\theta(y^{-} \mid x)}{\pi_{\mathrm{ref}}(y^{-} \mid x)}
  \right]
  \label{eq:logratio}
\end{equation}
and the loss
\begin{equation}
  \ell(\theta) = -\log \sigma(m_\theta)
    - \frac{\lambda}{|y^{+}|} \log \pi_\theta(y^{+} \mid x),
  \label{eq:loss}
\end{equation}
where $\log \pi(y \mid x)$ sums the log-probabilities of the tokens of
$y$, $\sigma$ is the logistic function and $|y^{+}|$ is the number of
tokens of $y^{+}$. Client $k$ minimizes the mean of $\ell$ over
$\mathcal{P}_k^{t}$ for \setting{pref.epochs} epoch, with $\beta$ set to
\setting{pref.beta} and $\lambda$ to \setting{pref.nll}, which yields
$\theta_k^{t}$.

The preference term is the DPO loss~\cite{rafailov2023dpo}, in which
$\beta$ sets the strength of the constraint toward the reference model.
We set $\pi_{\mathrm{ref}} = \pi_{\theta^{t-1}}$, the global model at
the start of the round. This reference model is identical for all
clients and is the model that generated the samples; as local training
starts from it, $m_\theta = 0$ for every pair initially. It is fixed
during the round, so its log-probabilities of the pairs are computed
once, before local training. The NLL term maintains the likelihood of
the correct sample, which the preference term constrains only relative
to that of the incorrect sample. Section~\ref{sec:ablation} ablates both
terms.

The values of $N$, the temperature, $\beta$ and $\lambda$, the uniform
pairing and the choice of reference model follow
CodeLutra~\cite{tao2024codelutra}, a centralized method. We differ in
constructing fewer pairs per question (CodeLutra constructs ten) and in
averaging the NLL term over tokens, where CodeLutra uses the
log-likelihood of the sequence.

Because sampling, labeling and pair construction are repeated in every
preference round, the pairs are on-policy: $\mathcal{P}_k^{t}$ reflects
the errors of $\pi_{\theta^{t-1}}$, whereas pairs constructed once would
reflect those of the model that generated them. Of the
\fact{cost.bird.pref1.questions} training questions of the BIRD clients
(Section~\ref{sec:setup}), \fact{cost.bird.pref1.paired} yield pairs in
the first preference round and \fact{cost.bird.pref2.paired} in the
second; for the Spider clients, \fact{cost.spider.pref1.paired} and
\fact{cost.spider.pref2.paired} of \fact{cost.spider.pref1.questions}.
On BIRD, the decrease stems mostly from questions on which all samples
are correct: their number grows from \fact{cost.bird.pref1.all_correct}
to \fact{cost.bird.pref2.all_correct}. Section~\ref{sec:ablation} also
ablates the on-policy pairs.

After local training, client $k$ uploads $\theta_k^{t}$ and its number
of pairs $|\mathcal{P}_k^{t}|$, and the server aggregates the adapters
as in FedAvg~\cite{mcmahan2017fedavg}:
\begin{equation}
  \theta^{t} = \sum_{k=1}^{K} w_k \, \theta_k^{t},
  \qquad
  w_k = \frac{|\mathcal{P}_k^{t}|}{\sum_{j=1}^{K} |\mathcal{P}_j^{t}|},
  \label{eq:aggregation}
\end{equation}
applied to every matrix of the adapter. FedAvg weights a client by its
amount of local data: $n_k$ in a warm-up round, the number of pairs
here. The samples, their execution results, the execution labels and the
preference pairs remain at the client; the size of the upload,
\setting{adapter.mb} for the adapter, equals that of a warm-up round.
